\documentclass[10pt,a4paper]{article}
\usepackage[utf8]{inputenc}
\usepackage[T1]{fontenc}
\usepackage[spanish]{babel}
\usepackage{geometry}
\geometry{left=1.5cm, right=1.5cm, top=1.5cm, bottom=1.5cm}
\usepackage{hyperref}
\usepackage{enumitem}
\usepackage{titlesec}

\titleformat{\section}{\large\bfseries\uppercase}{}{0em}{}[\titlerule]
\titlespacing{\section}{0pt}{1ex}{1ex}

\begin{document}
\pagestyle{empty}

\begin{center}
    {\huge \textbf{Jair Molina Arce}} \\
    \vspace{2mm}
    {\Large \textbf{Ingeniero Mecánico}} \\
    \vspace{2mm}
    Querétaro, Qro., México \\
    ingjairmolina@gmail.com | 5652646108 | jairmolina.dev | \href{https://github.com/smookymolina}{github.com/smookymolina} | \href{https://linkedin.com/in/jair-molina-arce-4909622b2}{linkedin.com/in/jair-molina-arce-4909622b2}
\end{center}

\section*{Perfil Profesional}
Ingeniero Mecánico egresado del IPN, con maestría en curso en Tecnologías Avanzadas y experiencia en diseño mecánico, análisis térmico y estructural (FEA), bancos de prueba e instrumentación. He caracterizado, diseñado y simulado sistemas mecánicos con SolidWorks, ANSYS y MATLAB/Simulink, y validado su comportamiento con adquisición de datos y pruebas experimentales. Mi diferencial es integrar la parte mecánica con electrónica y automatización: sensores, actuadores y microcontroladores trabajando como un solo sistema. Busco un rol de ingeniería mecánica en diseño, gestión térmica, pruebas o validación.

\section*{Educación}
\begin{itemize}[leftmargin=*, noitemsep]
    \item \textbf{Maestría en Tecnologías Avanzadas} (2024 -- Presente), Instituto Politécnico Nacional (IPN). Línea de investigación: control de sistemas dinámicos e instrumentación.
    \item \textbf{Ingeniería Mecánica} (2017 -- 2023), Instituto Politécnico Nacional (IPN). Enfoque en diseño de bancos de prueba, térmica, instrumentación y sistemas embebidos.
    \item \textbf{Bachillerato Técnico en Informática} (2014 -- 2017), Colegio de Bachilleres Plantel 1 ``El Rosario''.
\end{itemize}

\section*{Habilidades Técnicas}
\begin{itemize}[leftmargin=*, noitemsep]
    \item \textbf{Diseño / CAD-CAE:} SolidWorks, ANSYS (FEA), análisis estructural y térmico, MATLAB/Simulink, diseño y dimensionamiento de bancos de prueba, metrología, impresión 3D.
    \item \textbf{Térmica y fluidos:} Transferencia de calor, termodinámica aplicada, gestión térmica, control de ventilación, análisis energético.
    \item \textbf{Instrumentación / pruebas:} DAQ (LabVIEW, Python), sensores de presión/temperatura/flujo, NTC, control PID, osciloscopio, validación experimental.
    \item \textbf{Mecatrónica / automatización:} ESP32, STM32, RP2040/RP2350, C/C++, MicroPython, motores a pasos, MOSFET, UART, SPI, I2C, KiCad.
    \item \textbf{Datos y software:} Python (pandas, numpy, matplotlib), Excel avanzado (Power Query, VBA), Git, Docker.
\end{itemize}

\section*{Experiencia Relevante}
\begin{itemize}[leftmargin=*, noitemsep]
    \item \textbf{Docente - Ingeniería Térmica y Modelización de Sistemas Aeroespaciales} (Ago. 2025 -- Presente).
    \item \textbf{Docente - Tecnologías Disruptivas (Realidad Virtual y Aumentada)} (2024).
    \item \textbf{Ingeniería Mecánica e Instrumentación - Banco de Pruebas de Propulsión} (2022 -- 2024). FEA con ANSYS/SolidWorks, sistema DAQ.
\end{itemize}

\section*{Proyecto Mecánico Destacado}
\begin{itemize}[leftmargin=*, noitemsep]
    \item \textbf{Jaguar MX (2026):} Gestión térmica: dimensionamiento, control, esquemático y PCB. (Extracción de aire caliente, ventilador, NTC, FSM de seguridad, KiCad).
\end{itemize}

\section*{Proyectos Seleccionados}
\begin{itemize}[leftmargin=*, noitemsep]
    \item Pipeline de Automatización de Análisis de Datos y Reportes (2024 -- 2025).
    \item Aplicación de Realidad Virtual para Capacitación en Seguridad Industrial (2024).
    \item Sistema de Control Adaptativo PID con Ajuste Automático (2023 -- 2024).
    \item Impresora 3D DIY - Diseño y Construcción desde Cero.
\end{itemize}

\section*{Freelance}
\begin{itemize}[leftmargin=*, noitemsep]
    \item \textbf{Instructor} - Excel Avanzado (Mar. 2026).
    \item \textbf{Plataforma Web para Monitoreo Remoto de Variables Físicas} (2023 -- 2024).
\end{itemize}

\section*{Idiomas y Producción}
\begin{itemize}[leftmargin=*, noitemsep]
    \item \textbf{Idiomas:} Español nativo | Inglés avanzado con lectura técnica fluida.
    \item \textbf{Producción:} 75th IAC 2024 (DOI: 10.52202/078365-0120), Revista Hacia el Espacio (2024), CVU CONACYT \& ORCID.
\end{itemize}

\end{document}
